\documentclass[12pt]{article}

\usepackage{amssymb,amsmath,graphicx,subfig,float}
\usepackage[utf-8]{inputenc}
\usepackage[hidelinks]{hyperref}
\usepackage{geometry}
\geometry{left=1in, right=1in, top=1in, bottom=1in}
\usepackage{natbib}
\usepackage{xcolor}
\usepackage{pgfplots}
\usetikzlibrary{shapes.geometric, arrows.meta, positioning, fit}
\pgfplotsset{compat=1.18}
\usepackage{adjustbox}
\usepackage{booktabs}

% Unicode to LaTeX for math symbols
\usepackage{newunicodechar}
\newunicodechar{−}{\ensuremath{-}}
\newunicodechar{×}{\ensuremath{\times}}
\newunicodechar{→}{\ensuremath{\rightarrow}}
\newunicodechar{≈}{\ensuremath{\approx}}
\newunicodechar{≤}{\ensuremath{\leq}}
\newunicodechar{≥}{\ensuremath{\geq}}
\newunicodechar{±}{\ensuremath{\pm}}
\newunicodechar{α}{\ensuremath{\alpha}}
\newunicodechar{β}{\ensuremath{\beta}}
\newunicodechar{λ}{\ensuremath{\lambda}}
\newunicodechar{μ}{\ensuremath{\mu}}
\newunicodechar{σ}{\ensuremath{\sigma}}
\newunicodechar{≠}{\ensuremath{\neq}}

\title{Kinetic AI: Language Modeling as Equilibrium Computation}
\author{Sharath S. PhD}
\date{}

\begin{document}

\maketitle

\begin{abstract}
We propose Kinetic AI, a framework that reconceptualizes language model inference and training as solutions to equilibrium-seeking dynamical systems. Drawing on quantal response equilibrium theory, mechanism design, and deep equilibrium models, we ground language modeling in the language of game theory and economic theory. We implement the EqLM architecture, which computes equilibrium solutions implicitly via fixed-point iteration. We analyze convergence properties using maximum mean discrepancy, validate the approach on synthetic matrix games, and benchmark against standard pretraining regimes (BabyLM, GPT-2, BERT). This work opens a path to more interpretable, theoretically grounded language models.
\end{abstract}

\section{Introduction}

Language models predict the next token by computing a distribution over the vocabulary. We propose reinterpreting this computation as an agent solving an implicit equilibrium problem. Under a quantal response framework, the model is an economic agent with a utility function, and the distribution it learns is its best response to the data distribution it observes.

The key insight is that many learning dynamics can be understood as motion toward equilibrium in a game or mechanism where agents (or the model itself) optimize competing objectives. This view:
\begin{itemize}
\item Grounds inference in game-theoretic principles, making objectives interpretable.
\item Connects training to auction theory and incentive design.
\item Enables convergence analysis via mathematical tools not typically applied to neural networks.
\item Suggests new architectural designs that enforce equilibrium properties by construction.
\end{itemize}

In this paper, we develop this intuition into a concrete architecture and set of experiments.

\section{Background}

\subsection{Quantal Response Equilibrium}

Quantal response equilibrium (QRE) is a refinement of Nash equilibrium that models bounded rationality. Instead of best-response, agents play a strategy that is a noisy function of their payoff: the probability of playing action $a$ is proportional to $\exp(\beta u(a))$, where $u(a)$ is the utility and $\beta$ is an inverse temperature parameter.

In language modeling, we can view the model as a noisy best responder to the data distribution, with the logit layer implementing a softmax that is precisely this quantal response.

\subsection{Maximum Mean Discrepancy}

Maximum mean discrepancy (MMD) is a kernel-based divergence between distributions. For distributions $P$ and $Q$ over $\mathcal{X}$:
\[
\text{MMD}(P, Q) = \left\lVert \mathbb{E}_{x \sim P}[\phi(x)] - \mathbb{E}_{x \sim Q}[\phi(x)] \right\rVert_{\mathcal{H}}
\]
where $\phi$ is a feature map into a reproducing kernel Hilbert space. MMD is zero if and only if $P = Q$.

We use MMD to analyze convergence of learned distributions toward target equilibrium distributions.

\subsection{Deep Equilibrium Models}

Deep Equilibrium (DEQ) models replace explicit depth with implicit layers: a single function is iterated to convergence via fixed-point solvers. Given a fixed map $f_\theta(\mathbf{z}, \mathbf{x})$, the equilibrium is defined as $\mathbf{z}^* = f_\theta(\mathbf{z}^*, \mathbf{x})$. DEQ models achieve strong performance on vision and NLP tasks while offering interpretability: the model solves an implicit problem.

We adopt the DEQ perspective: language model layers can be viewed as steps toward an implicit equilibrium.

\subsection{Mechanism Design for Language Models}

Recent work (Duetting et al. 2024) frames language model alignment as a mechanism design problem: the model must be incentivized to produce outputs that align with human preferences. We draw on this framework to motivate architectural constraints that enforce certain equilibrium properties.

\section{The EqLM Architecture}

The Equilibrium Language Model (EqLM) is a Transformer-based architecture augmented with an implicit equilibrium layer.

\subsection{Core Design}

The model processes input tokens through standard Transformer blocks to obtain contextual representations $\mathbf{h}$. At the final layer, instead of a single linear projection to logits, we solve:
\[
\mathbf{z}^* = f(\mathbf{z}, \mathbf{h}; \theta)
\]
for $\mathbf{z}^* \in \mathbb{R}^V$ (vocabulary dimension), where $f$ encodes game-theoretic constraints (e.g., the softmax normalization ensures QRE structure). The output distribution is $\pi = \text{softmax}(\mathbf{z}^*)$.

\subsection{Fixed-Point Iteration}

We compute $\mathbf{z}^*$ via Anderson acceleration or Newton-Schulz iteration, maintaining interpretability: the model is not a ``black box'' but a solution to an explicit optimization problem.

\section{Convergence Analysis (Maximum Mean Discrepancy)}

We analyze convergence of the learned token distribution $\pi_\text{model}$ to a target distribution $\pi_\text{data}$ using MMD as a metric.

% Filled from validated findings at phase closure — see research/memory/findings.md. No numbers without a real run.

\section{Experiments}

\subsection{Matrix Game Convergence}

% Filled from validated findings at phase closure — see research/memory/findings.md. No numbers without a real run.

\subsection{Deep Equilibrium Memory and Computational Cost}

% Filled from validated findings at phase closure — see research/memory/findings.md. No numbers without a real run.

\subsection{BabyLM Pretraining vs.\ GPT-2 and BERT}

% Filled from validated findings at phase closure — see research/memory/findings.md. No numbers without a real run.

\subsection{Auction-Based Decoding: MPO vs.\ DPO}

% Filled from validated findings at phase closure — see research/memory/findings.md. No numbers without a real run.

\section{Related Work}

The connection between language modeling and equilibrium is implicit in much recent work. Sokota et al.\ (2023) provides a unified framework connecting RL, QRE, and zero-sum games. Bai, Kolter, and Koltun (2019) pioneered DEQ models. Duetting et al.\ (2024) frame LLM alignment as mechanism design. Direct Preference Optimization (Rafailov et al., 2023) and Self-Play Preference Optimization (Wu et al., 2024) can be reinterpreted as equilibrium-seeking in preference space.

We also build on classical results: McKelvey and Palfrey (1995) on quantal response equilibria, Anderson (1965) on iterative procedures, and Hoda et al.\ (2010) on smoothing in sequential games.

\section{Limitations}

\begin{itemize}
\item Convergence to fixed points is computationally expensive; practical inference requires approximate solvers.
\item The framework assumes a well-defined game structure, which may not hold for all language modeling tasks.
\item We do not yet achieve state-of-the-art performance on standard benchmarks; the focus is interpretability and theory.
\end{itemize}

\section{Conclusion}

We have proposed Kinetic AI, a framework for language modeling grounded in game theory and equilibrium computation. The EqLM architecture makes this concrete, and preliminary experiments validate the approach. This work opens new directions for interpretable, theoretically-grounded language models.

Future work includes scaling EqLM to larger models, applying the framework to alignment and preference learning, and exploring connections to other dynamical systems perspectives on neural networks.

\bibliographystyle{plainnat}
\bibliography{references}

\end{document}
